\section{Introduction}
The Kitaev honeycomb model~\cite{Kitaev} is an exactly solvable model that hosts a quantum spin liquid ground state~\cite{Savary_2017, RMP-QSL}, which possesses features of topological order, such as anyon statistics~\cite{Kitaev, JW-2, anyon-2, anyon-3}, topological ground state degeneracy~\cite{degeneracy-1, degeneracy-2, Mandal-Shankar-Baskaran-2012}, and fractionalization of degrees of freedom~\cite{JW-1, fractionalization-4}, which is even believed to be observed in experiments~\cite{Kitaev-review, fractionalization-1, fractionalization-2, fractionalization-3}. Numerous studies have solved the Kitaev model by transforming spin operators to fermions, including parton construction and Jordan-Wigner transform~\cite{burnell-nayak, feng-zhang-xiang, JW-1, JW-2, Fu-Knolle-Perkins-2018}. A quantum spin liquid state is a phase of matter in which the standard Ginzburg-Landau paradigm breaks down since no symmetry is broken. Therefore, a theoretical framework beyond the Ginzburg-Landau paradigm is necessary to describe this novel phase of matter, and quantum entanglement serves as an insightful diagnostic of its underlying structure~\cite{QuantumOrderEntanglementPattern-1, QuantumOrderEntanglementPattern-2, QuantumOrderEntanglementPattern-3, QuantumOrderEntanglementPattern-4}. 

Many-body entanglement and its application to modern condensed matter physics have been extensively studied in the last two decades~\cite{entanglement-in-many-body, area-law-for-EE, quantum-entanglement-in-condensed-matter}. In particular, it is crucial in the quantum topological phases of matter~\cite{RMP-Zoo}. Several quantities on entanglement have been proven to be related to features of quantum matters, such as entanglement entropy, in which some specific terms encode information of topological order~\cite{topoEE_1, topoEE_2, FlammiaHammaAlioscia-topoEE-renyi-flat-projector}, and the entanglement spectrum, whose low-lying portion is conjectured to hint at bulk-boundary correspondence~\cite{Li-Haldane, Li-Haldane-1, Li-Haldane-2}. These studies of entanglement stimulated a broader research field that applies a quantum information theory approach to investigate the structures of quantum phases of matter~\cite{ZengChenZhouWen2019}. In this framework, a more convenient way to represent a state would be a density matrix in an operator space instead of a vector in a Hilbert space, not only because of its ability to describe a mixed state, but also the convenience it brings for analysis using operator algebras~\cite{BratteliRobinson1987, Naaijkens2017}. More importantly, the density matrix provides a concise and efficient language for the study of entanglement. Since the entanglement is characterized by correlations between a subsystem and another one, the reduced density matrix is the only object capable of encoding the state of the subsystem. Therefore, investigating the density matrix and the information encoded therein is essential to understanding the entanglement structure of a state.

Among various models, the entanglement properties of free fermion and free boson gas have been extensively studied, including the reduced density matrix and the entanglement spectrum~\cite{Peschel_2003, Peschel_2009, Chung_2001, Cheong_2004}, as well as the entanglement entropy~\cite{EE-free-boson-1, EE-free-boson-2, EE-free-fermion-1, EE-free-fermion-2}. Due to this, the many-body entanglement of certain spin-1/2 systems can also be studied in the same manner~\cite{Jin-Korepin, Vidal-Latorre-Rico-Kitaev, Latorre-Rico-Vidal}, since there exists a duality between a spin system of that type and a free-fermion system. In light of this, the entanglement structure of the Kitaev spin liquid is supposed to be able to be transformed to that of a free fermion system because the Kitaev model is equivalent to a free fermion system in the ground state. The discussion of the entanglement entropy of the Kitaev model in this sense has been given by \cite{Yao-Qi}. The entanglement entropy is shown to be separated into the gauge part and the fermion part, in which the fermion part can be analyzed using the treatment for free fermion systems. In addition, the scaling of the entanglement entropy and the entanglement spectrum of the equivalent fermion model have also been investigated~\cite{Kitaev-entanglement-1, Kitaev-entanglement-2}.

Although substantial progress has been made in understanding the physical and entanglement properties of the Kitaev model, several important issues remain to be addressed. First, the majority of the existing literature analyzes the Kitaev model in the fermion representation. However, how the entanglement properties are reflected by spin degrees of freedom--the true constituents of the system--remains unclear. Second, although~\cite{Yao-Qi} has given a comprehensive picture of the entanglement feature of the fermion sector of the Kitaev ground state, it should be noted that the approach adopted there does not directly establish the equivalence between the entanglement spectra of the spin system and the corresponding fermion system. Moreover, the derivation of the separability provided by~\cite{Yao-Qi} relies on the assumption that the subregion contains an even number of lattice sites, leaving the odd-site case unexplored. Third, the contribution to the entanglement from the gauge field is not investigated as extensively as that from the fermion sector. This problem has been pointed out by~\cite{Dora-Moessner-2018}, and the entanglement of the gauge field is discussed in the paper based on the Jordan-Wigner formulation. Nevertheless, the analysis still relies on the fermionic formulation, and the origin and physical interpretation of gauge-field entanglement are not transparent within the spin-language description. Lastly, a detailed discussion of the density matrix of the Kitaev ground state, especially its explicit expression in spin language, remains absent. Several works have discussed the density matrix of the Kitaev ground state using spin language, but only concentrated on the two-point reduced density matrix~\cite{correlation-densitymatrix-1, correlation-densitymatrix-2}. However, a two-point correlation function is not sufficient to detect topological order because topological information is usually hidden in loop structures and anyon statistics~\cite{toric-code, string-net, topoEE_1, topoEE_2, Hasting-Wen, chiral-spin-liquid}. 

In this work, we thoroughly study the density matrix of the Kitaev ground state and its entanglement properties using spin language. In Sections~\ref{sec:density-matrix-whole}, we introduce two important concepts in this work: (i) equivalence classes and (ii) string operators. We find that the density matrix is constructed from string operators, and the equivalence relations among these operators simplify its final form. In \ref{sec:density-matrix-calculate}, we derive the density matrix of the Kitaev ground state on a torus. Afterward, the reduced density matrix of a simply-connected subregion is also obtained in Section~\ref{sec:density-matrix-reduced}. We will see in Section~\ref{sec:Emergent-Fermions-at-the-Endpoints} that the form of the density matrix shows an evident correspondence between the Majorana fermion representation and the spin representation of the Kitaev model, which reflects the physical picture in which the fermions emerge at endpoints of strings \cite{fermion-and-string, string-net, Photons-and-electrons-as-emergent-phenomena-string, PRD-string-wen-quantum-order}. In Section~\ref{sec:EE-calculation-even}, we prove that the entanglement entropy obtained by directly computing the von Neumann entropy of a reduced density matrix reproduces the result obtained in~\cite{Yao-Qi}. Moreover, in Section~\ref{sec:EE-calculation-odd}, our method is shown to be valid even for subregions with an odd number of sites. We also point out an additional 2-fold degeneracy in the entanglement spectrum provided by the fermion sector, and it turns out that the fermion parity plays a crucial role. In Section~\ref{sec:origin-of-degeneracy}, we prove that the separability of the entanglement entropy is due to the block-diagonal structure of the density matrix and the extensive degeneracy in the entanglement spectrum, which is discussed along the lines of the mechanism proposed by~\cite{PRX-ent-high}, which states the exact 1-form Wilson symmetry together with the exact Gauss law of the gauge theory gives this extensive degeneracy. In the same section, the physical interpretation of why fermion parity is crucial in cases with an odd number of lattice sites is elucidated using the spin formulation.


